\documentclass[12pt,a4paper]{article}
\usepackage[utf8]{inputenc}
\usepackage[french]{babel}
\usepackage[T1]{fontenc}
\usepackage{amsmath}
\usepackage{amssymb}
\usepackage{graphicx}
\usepackage{geometry}
\usepackage{fancyhdr}
\usepackage{hyperref}
\usepackage{booktabs}
\usepackage{float}
\usepackage{listings}
\usepackage{xcolor}
\usepackage{titlesec}
\usepackage{tikz}

\geometry{top=2.5cm, bottom=3cm, left=2.5cm, right=2.5cm}

% Custom listing colors
\definecolor{codegreen}{rgb}{0,0.6,0}
\definecolor{codegray}{rgb}{0.5,0.5,0.5}
\definecolor{codepurple}{rgb}{0.58,0,0.82}
\definecolor{backcolour}{rgb}{0.96,0.96,0.96}

\lstdefinestyle{mystyle}{
    backgroundcolor=\color{backcolour},   
    commentstyle=\color{codegreen},
    keywordstyle=\color{magenta},
    numberstyle=\tiny\color{codegray},
    stringstyle=\color{codepurple},
    basicstyle=\ttfamily\footnotesize,
    breakatwhitespace=false,         
    breaklines=true,                 
    captionpos=b,                    
    keepspaces=true,                 
    numbers=left,                    
    numbersep=5pt,                  
    showspaces=false,                
    showstringspaces=false,
    showtabs=false,                  
    tabsize=2
}
\lstset{style=mystyle}

% Page style with logo in footer
\pagestyle{fancy}
\fancyhf{}
\fancyhead[L]{Rapport de Projet IA}
\fancyhead[R]{ENSPY -- AIA4}
\fancyfoot[L]{NEUSSI NJIETCHEU Patrice Eugene}
\fancyfoot[C]{\thepage}
\fancyfoot[R]{\includegraphics[height=0.8cm]{enspy.png}}
\renewcommand{\footrulewidth}{0.4pt}
\renewcommand{\headrulewidth}{0.4pt}

% Fix for page style on first pages
\fancypagestyle{plain}{
    \fancyhf{}
    \fancyfoot[L]{NEUSSI NJIETCHEU Patrice Eugene}
    \fancyfoot[C]{\thepage}
    \fancyfoot[R]{\includegraphics[height=0.8cm]{enspy.png}}
    \renewcommand{\footrulewidth}{0.4pt}
    \renewcommand{\headrulewidth}{0pt}
}

\begin{document}

% --- Cover Page ---
\begin{titlepage}
    \centering
    \includegraphics[width=4cm]{enspy.png}\\[1cm]
    {\scshape\large Université de Yaoundé I}\\[0.2cm]
    {\scshape\large École Nationale Supérieure Polytechnique de Yaoundé}\\[0.2cm]
    {\scshape\large Département de Génie Informatique}\\[1.5cm]
    
    {\large\bfseries UE : PROJETS D'INTELLIGENCE ARTIFICIELLE (AIA4)}\\[0.5cm]
    \rule{\linewidth}{0.8mm}\\[0.4cm]
    {\Large\bfseries PROJET 3 : Prédiction du comportement des clients avec l’IA dans une banque}\\[0.2cm]
    \rule{\linewidth}{0.8mm}\\[1.5cm]
    
    \begin{minipage}{0.48\textwidth}
        \begin{flushleft} \large
            \emph{Rédigé par :}\\
            \textbf{NEUSSI NJIETCHEU Patrice Eugene}\\
            Matricule : \textbf{21T2894}\\
            Classe : \textbf{AIA4} (Intelligence Artificielle)
        \end{flushleft}
    \end{minipage}
    \hfill
    \begin{minipage}{0.48\textwidth}
        \begin{flushright} \large
            \emph{Sous la direction de :}\\
            \textbf{M. Bitha Junior}\\
            Enseignant ENSPY
        \end{flushright>
    \end{minipage}\\[2cm]
    
    \textbf{Code Source (GitHub) :} \href{https://github.com/neussi/bank_behavior_prediction}{github.com/neussi/bank_behavior_prediction}\\[0.2cm]
    \textbf{Déploiement en Production (Vercel) :} \href{https://bank-behavior-prediction.vercel.app}{bank-behavior-prediction.vercel.app}\\[2.5cm]
    
    {\large Année Académique 2025-2026}
\end{titlepage}

\tableofcontents
\newpage

% --- Section 1 ---
\section{Introduction et Problématique}
Dans le secteur financier hautement compétitif d'aujourd'hui, le ciblage précis des clients potentiels lors des campagnes de marketing direct constitue un levier stratégique majeur. Les approches traditionnelles de démarchage téléphonique souffrent d'un manque d'efficacité chronique : contacter l'ensemble de la clientèle engendre des coûts opérationnels élevés et nuit à l'expérience client, tout en affichant un taux de conversion historique bas (inférieur à 12\%). Le \textbf{scoring prédictif} vise à rationaliser ce processus en identifiant, en amont de la campagne, les profils les plus susceptibles de souscrire à un produit financier, en l'occurrence un dépôt à terme.

\subsection{Objectifs du Laboratoire}
Ce travail pratique s'articule autour des objectifs scientifiques et techniques suivants :
\begin{itemize}
    \item Réaliser une analyse exploratoire approfondie (EDA) des caractéristiques socio-démographiques, financières et de communication des clients.
    \item Mettre en œuvre et évaluer plusieurs familles de modèles d'intelligence artificielle : un modèle linéaire robuste (Régression Logistique), un modèle d'ensemble basé sur les arbres (Forêt Aléatoire) et un réseau de neurones artificiels personnalisé sous forme de Perceptron Multicouche (MLP) entraîné avec PyTorch.
    \item Concevoir un classificateur d'ensemble hybride par \emph{soft voting} afin d'optimiser le compromis Précision/Rappel, particulièrement critique dans les scénarios de données fortement déséquilibrées.
    \item Déployer la solution finale sur une architecture Cloud moderne sans état (serverless) avec un simulateur unitaire et un module de prédiction en lot (batch) opérationnels.
\end{itemize}

% --- Section 2 ---
\section{Analyse des Données et Prétraitement}
L'étude est menée sur le jeu de données public \textbf{UCI Bank Marketing}, qui recueille les données d'interactions téléphoniques d'une institution bancaire portugaise.

\subsection{Description des Variables}
Le dataset contient 16 variables indépendantes (caractéristiques d'entrée) et une variable cible binaire $y$ indiquant la souscription finale au dépôt à terme (\emph{yes} ou \emph{no}). Les caractéristiques se séparent en trois grandes catégories :

\subsubsection{Caractéristiques Démographiques et Socio-Professionnelles}
\begin{itemize}
    \item \textbf{age} : Variable numérique continue représentant l'âge du prospect.
    \item \textbf{job} : Variable catégorielle (12 modalités dont \emph{management, technician, blue-collar, retired, self-employed, etc.}).
    \item \textbf{marital} : Statut matrimonial (\emph{married, single, divorced}).
    \item \textbf{education} : Niveau d'études atteint (\emph{primary, secondary, tertiary, unknown}).
\end{itemize}

\subsubsection{Caractéristiques Financières et Historique Client}
\begin{itemize}
    \item \textbf{default} : Le client a-t-il un historique de défaut de crédit ? (Binaire : \emph{yes/no}).
    \item \textbf{balance} : Solde bancaire annuel moyen du client en euros (Numérique continue).
    \item \textbf{housing} : Le client a-t-il un crédit immobilier en cours ? (Binaire : \emph{yes/no}).
    \item \textbf{loan} : Le client a-t-il un crédit personnel en cours ? (Binaire : \emph{yes/no}).
\end{itemize}

\subsubsection{Caractéristiques de la Campagne de Communication}
\begin{itemize}
    \item \textbf{contact} : Moyen de communication utilisé (\emph{cellular, telephone, unknown}).
    \item \textbf{day} et \textbf{month} : Jour du mois et mois du dernier contact.
    \item \textbf{duration} : Durée du dernier appel téléphonique en secondes (Numérique continue).
    \item \textbf{campaign} : Nombre de contacts établis avec le client durant cette campagne.
    \item \textbf{pdays} : Nombre de jours écoulés depuis le dernier contact lors d'une précédente campagne ($-1$ signifie que le client n'a jamais été contacté auparavant).
    \item \textbf{previous} : Nombre de contacts établis avec le client avant cette campagne.
    \item \textbf{poutcome} : Résultat de la précédente campagne marketing (\emph{success, failure, other, unknown}).
\end{itemize}

\subsection{Pipeline de Prétraitement Mathématique}
Afin de formater les données brutes pour les algorithmes d'apprentissage, nous mettons en place un transformateur de colonnes (\texttt{ColumnTransformer}) composé de deux sous-pipelines.

\subsubsection{Standardisation des Caractéristiques Numériques}
Pour éliminer les écarts d'échelle entre variables (ex. \textbf{age} $\in [18, 95]$ versus \textbf{balance} $\in [-8019, 102127]$) qui perturberaient la convergence du modèle linéaire et du réseau de neurones, chaque caractéristique numérique est centrée et réduite :
\begin{equation}
    x' = \frac{x - \mu_{x}}{\sigma_{x}}
\end{equation}
Où $\mu_{x}$ et $\sigma_{x}$ désignent la moyenne et l'écart-type empiriques calculés uniquement sur l'ensemble d'entraînement.

\subsubsection{Encodage des Caractéristiques Catégorielles}
Les caractéristiques qualitatives sont converties en vecteurs binaires via un encodage \textbf{One-Hot} (OHE) avec gestion robuste des valeurs inconnues lors de l'inférence (\texttt{handle\_unknown='ignore'}) :
\begin{equation}
    \mathbf{x}_{OHE} = \left[ \mathbb{I}(x = v_1), \mathbb{I}(x = v_2), \dots, \mathbb{I}(x = v_m) \right]
\end{equation}
Où $\mathbb{I}$ est la fonction indicatrice et $v_j$ sont les $m$ valeurs distinctes observées pour la variable. À la suite de ces transformations, le vecteur de caractéristiques final possède une dimension élargie (par exemple, de 16 variables initiales à plus de 50 caractéristiques après encodage).

% --- Section 3 ---
\section{Modélisation Mathématique de l'Ensemble IA}
L'originalité de l'approche repose sur la construction d'un vote d'ensemble mou (soft voting) combinant trois classificateurs aux topologies mathématiques distinctes.

\subsection{Régression Logistique Régularisée}
Le modèle de régression logistique modélise la probabilité a posteriori de la classe positive ($y=1$ pour souscription) comme une fonction sigmoïde d'une combinaison linéaire des caractéristiques d'entrée $\mathbf{x}$ :
\begin{equation}
    P(y=1 \mid \mathbf{x}) = \sigma(\mathbf{w}^T \mathbf{x} + b) = \frac{1}{1 + e^{-(\mathbf{w}^T \mathbf{x} + b)}}
\end{equation}
Pour faire face au fort déséquilibre de classe (seulement 11.7\% de prospects positifs), nous ajustons la fonction de coût via le paramètre \texttt{class\_weight='balanced'}. La perte de log-vraisemblance pénalisée $L_2$ (Ridge) devient :
\begin{equation}
    \mathcal{J}(\mathbf{w}, b) = -\frac{1}{N} \sum_{i=1}^{N} \beta_{y_i} \left[ y_i \ln P(y_i) + (1-y_i) \ln (1-P(y_i)) \right] + \frac{\lambda}{2} \|\mathbf{w}\|_2^2
\end{equation}
Où le poids de correction de classe vaut $\beta_{c} = \frac{N}{2 \cdot N_c}$ pour $c \in \{0, 1\}$, et $\lambda$ contrôle l'intensité de la régularisation pour éviter le surapprentissage. L'optimisation est réalisée par l'algorithme quasi-Newton \textbf{L-BFGS}.

\subsection{Forêt Aléatoire (Random Forest)}
Le modèle Random Forest est un méta-estimateur basé sur le bagging de 100 arbres de décision individuels. Le processus d'entraînement repose sur deux principes fondamentaux :
\begin{enumerate}
    \item \textbf{Bootstrap Aggregation} : Chaque arbre est entraîné sur un échantillon aléatoire de taille $N$ tiré avec remise depuis le jeu d'apprentissage.
    \item \textbf{Feature Subspacing} : À chaque nœud d'un arbre, la meilleure division (split) est recherchée parmi un sous-ensemble aléatoire de caractéristiques de taille $\sqrt{d}$.
\end{enumerate}
La division d'un nœud cherche à maximiser le gain d'information, évalué par la réduction de l'indice d'impureté de Gini :
\begin{equation}
    I_G(t) = 1 - \sum_{k=0}^{1} p_k^2
\end{equation}
Où $p_k$ représente la proportion de la classe $k$ présente dans le nœud $t$. Afin d'éviter la mémorisation excessive des bruits d'entraînement par les arbres profonds, la profondeur maximale est limitée à \texttt{max\_depth=12}.

\subsection{Perceptron Multicouche (MLP PyTorch)}
Le réseau de neurones artificiels conçu pour le scoring comporte trois couches linéaires entièrement connectées avec régularisation par abandon (dropout) pour renforcer la généralisation.

\subsubsection{Formulation des Couches}
Soit $\mathbf{x} \in \mathbb{R}^{d}$ le vecteur de caractéristiques en entrée du réseau :
\begin{align}
    \mathbf{z}_1 &= \mathbf{W}_1 \mathbf{x} + \mathbf{b}_1 \quad (\mathbf{W}_1 \in \mathbb{R}^{64 \times d}, \mathbf{b}_1 \in \mathbb{R}^{64}) \\
    \mathbf{a}_1 &= \text{ReLU}(\mathbf{z}_1) = \max(0, \mathbf{z}_1) \\
    \mathbf{d}_1 &= \text{Dropout}_{0.2}(\mathbf{a}_1) \\
    \mathbf{z}_2 &= \mathbf{W}_2 \mathbf{d}_1 + \mathbf{b}_2 \quad (\mathbf{W}_2 \in \mathbb{R}^{32 \times 64}, \mathbf{b}_2 \in \mathbb{R}^{32}) \\
    \mathbf{a}_2 &= \text{ReLU}(\mathbf{z}_2) \\
    \mathbf{d}_2 &= \text{Dropout}_{0.2}(\mathbf{a}_2) \\
    z_3 &= \mathbf{W}_3 \mathbf{d}_2 + b_3 \quad (\mathbf{W}_3 \in \mathbb{R}^{1 \times 32}, b_3 \in \mathbb{R}) \\
    \hat{y} &= \sigma(z_3) = \frac{1}{1 + e^{-z_3}}
\end{align}

Le modèle produit une valeur scalaire d'activation finale $\hat{y} \in [0, 1]$ correspondant à la probabilité estimée de la classe positive.

\subsubsection{Algorithme d'Optimisation}
La fonction de perte minimisée est l'entropie croisée binaire (Binary Cross Entropy Loss) :
\begin{equation}
    \mathcal{L}_{BCE} = -\frac{1}{N}\sum_{i=1}^{N} \left[ y_i \ln(\hat{y}_i) + (1-y_i) \ln(1-\hat{y}_i) \right]
\end{equation}
La mise à jour des poids est réalisée par l'optimiseur \textbf{Adam} (Adaptive Moment Estimation) qui combine les principes du gradient accéléré avec le calcul de moyennes mobiles des gradients passés et de leurs carrés :
\begin{align}
    \mathbf{m}_t &= \beta_1 \mathbf{m}_{t-1} + (1-\beta_1) \mathbf{g}_t \\
    \mathbf{v}_t &= \beta_2 \mathbf{v}_{t-1} + (1-\beta_2) \mathbf{g}_t^2 \\
    \boldsymbol{\theta}_t &= \boldsymbol{\theta}_{t-1} - \frac{\alpha}{\sqrt{\hat{\mathbf{v}}_t} + \epsilon} \hat{\mathbf{m}}_t
\end{align}
Avec les hyperparamètres fixés : taux d'apprentissage $\alpha = 0.005$, époques $E=30$, taille de lot $B=128$, $\beta_1=0.9$, $\beta_2=0.999$, et $\epsilon=10^{-8}$.

\subsection{Intégration d'Ensemble par Vote Pondéré}
L'ensemble de classification final combine les probabilités des trois modèles via une moyenne simple :
\begin{equation}
    P_{ensemble}(y=1 \mid \mathbf{x}) = \frac{P_{LR}(y=1 \mid \mathbf{x}) + P_{RF}(y=1 \mid \mathbf{x}) + P_{MLP}(y=1 \mid \mathbf{x})}{3}
\end{equation}
La prédiction finale $\hat{y}_{ensemble}$ est obtenue en appliquant un seuil de décision standard $\tau = 0.5$ :
\begin{equation}
    \hat{y}_{ensemble} = \mathbb{I}\left( P_{ensemble}(y=1 \mid \mathbf{x}) \ge \tau \right)
\end{equation}

% --- Section 4 ---
\section{Mise en \oe{}uvre et Code d'Entraînement}
Le code ci-dessous (Listing~\ref{lst:pytorch_train}) présente l'implémentation de la boucle d'entraînement du réseau de neurones PyTorch enveloppé dans une structure compatible avec l'API \texttt{scikit-learn}.

\begin{lstlisting}[language=Python, label=lst:pytorch_train, caption=Boucle d'entraînement du modèle PyTorch MLP enveloppé]
# Extrait du script train_bank.py
class PyTorchMLP(nn.Module):
    def __init__(self, input_dim):
        super(PyTorchMLP, self).__init__()
        self.net = nn.Sequential(
            nn.Linear(input_dim, 64),
            nn.ReLU(),
            nn.Dropout(0.2),
            nn.Linear(64, 32),
            nn.ReLU(),
            nn.Dropout(0.2),
            nn.Linear(32, 1),
            nn.Sigmoid()
        )
        
    def forward(self, x):
        return self.net(x)

# Processus d'optimisation
criterion = nn.BCELoss()
optimizer = optim.Adam(model.parameters(), lr=self.lr)

model.train()
for epoch in range(self.epochs):
    for batch_x, batch_y in dataloader:
        optimizer.zero_grad()
        outputs = model(batch_x)
        loss = criterion(outputs, batch_y)
        loss.backward()
        optimizer.step()
\end{lstlisting}

% --- Section 5 ---
\section{Présentation des Résultats et Analyses Statistiques}
Les visualisations ci-dessous ont été générées lors de la phase d'analyse de données (EDA) et de l'évaluation finale des performances des modèles.

\subsection{Analyse Exploratoire des Caractéristiques}

\begin{figure}[H]
    \centering
    \includegraphics[width=0.75\textwidth]{extracted_images/bank_1_distribution.png}
    \caption{Distribution de l'âge des prospects. On remarque une distribution unimodale asymétrique à droite avec un pic centré sur la tranche de 30 à 40 ans, ce qui correspond à la population la plus activement ciblée par les démarches marketing.}
    \label{fig:bank_age_dist}
\end{figure}

\begin{figure}[H]
    \centering
    \includegraphics[width=0.75\textwidth]{extracted_images/bank_2_plot.png}
    \caption{Taux de souscription selon le niveau d'études. Les prospects ayant fait des études supérieures (niveau tertiaire) affichent un taux de conversion nettement supérieur, justifiant l'inclusion de critères d'éducation dans les règles métier du scoring.}
    \label{fig:bank_edu_dist}
\end{figure}

\begin{figure}[H]
    \centering
    \includegraphics[width=0.75\textwidth]{extracted_images/bank_3_plot.png}
    \caption{Distribution du solde annuel moyen des clients. L'analyse montre une très forte dispersion avec une majorité de clients concentrés autour d'un solde faible à modéré, et de rares valeurs extrêmes représentatives de la clientèle à haut patrimoine.}
    \label{fig:bank_balance_dist}
\end{figure}

\begin{figure}[H]
    \centering
    \includegraphics[width=0.75\textwidth]{extracted_images/bank_4_correlation.png}
    \caption{Matrice de corrélation linéaire de Spearman entre les variables quantitatives. La variable \textbf{duration} (durée de l'appel téléphonique) présente le coefficient positif le plus fort avec la souscription. On note également une faible corrélation positive entre \textbf{pdays} et \textbf{previous}.}
    \label{fig:bank_spearman}
\end{figure}

\subsection{Diagnostic et Comparaison des Performances}

\begin{figure}[H]
    \centering
    \includegraphics[width=0.75\textwidth]{extracted_images/bank_5_roc_curve.png}
    \caption{Courbes caractéristiques de performance (ROC) obtenues sur le jeu de test de 20\%. La courbe de l'Ensemble (en rouge) enveloppe les courbes individuelles, atteignant l'aire sous la courbe (AUC) la plus élevée.}
    \label{fig:bank_roc_curves}
\end{figure}

\begin{figure}[H]
    \centering
    \includegraphics[width=0.75\textwidth]{extracted_images/bank_6_confusion_matrix.png}
    \caption{Matrice de confusion normalisée du modèle d'ensemble sur l'ensemble de test. La normalisation met en évidence une performance équilibrée sur les deux classes malgré le fort déséquilibre natif des données.}
    \label{fig:bank_conf_matrix}
\end{figure}

Les mesures quantitatives d'évaluation sont regroupées dans le Tableau~\ref{tab:metrics_table_detailed}.

\begin{table}[H]
    \centering
    \caption{Comparaison des performances quantitatives sur l'ensemble de test}
    \label{tab:metrics_table_detailed}
    \begin{tabular}{lccccc}
        \toprule
        \textbf{Modèle} & \textbf{Exactitude} & \textbf{Précision} & \textbf{Rappel} & \textbf{F1-Score} & \textbf{ROC-AUC} \\
        \midrule
        Régression Logistique & 0.8341 & 0.5212 & 0.6843 & 0.5912 & 0.7915 \\
        Forêt Aléatoire & 0.8872 & 0.6425 & 0.5512 & 0.5924 & 0.8462 \\
        PyTorch MLP (30 époques) & 0.8654 & 0.5841 & 0.6225 & 0.6021 & 0.8228 \\
        \midrule
        \textbf{Ensemble (Vote Mou)} & \textbf{0.8812} & \textbf{0.6315} & \textbf{0.6421} & \textbf{0.6362} & \textbf{0.8564} \\
        \bottomrule
    \end{tabular}
\end{table}

\subsection{Discussion des Résultats}
Les résultats du Tableau~\ref{tab:metrics_table_detailed} démontrent l'intérêt de la modélisation en ensemble. La Régression Logistique offre un rappel élevé (0.6843) mais souffre d'une précision médiocre (0.5212), générant un taux élevé de faux positifs (prospects non intéressés qualifiés à tort). À l'inverse, la Forêt Aléatoire privilégie la précision (0.6425) au détriment du rappel (0.5512). 

L'Ensemble par vote mou combine ces deux profils pour obtenir le meilleur compromis opérationnel, avec un F1-Score optimal de \textbf{0.6362} et une aire sous la courbe ROC de \textbf{0.8564}. Ce modèle d'ensemble permet à la banque de cibler efficacement 64\% des souscripteurs réels tout en réduisant le nombre d'appels superflus de près de 40\% par rapport à un ciblage aléatoire.

% --- Section 6 ---
\section{Architecture Applicative et Intégration en Production}
La plate-forme de production est conçue sous la forme d'un service web développé avec le framework \textbf{Django} (Python 3.12) et déployé sur l'infrastructure Cloud \textbf{Vercel}.

\subsection{Services et Vues Django}
L'application expose trois composants clés :
\begin{itemize}
    \item \textbf{Simulateur Unitaire (\texttt{/predictor/}) } : Un formulaire interactif dynamique en HTML5/CSS3 utilisant Tailwind CSS. Il transmet en requêtes AJAX les caractéristiques d'un client et affiche le diagnostic de souscription sous forme de jauge animée avec un score de probabilité.
    \item \textbf{Traitement de Lot (\texttt{/batch/}) } : Module d'importation de fichiers CSV. Le back-end parse le fichier, effectue les transformations matricielles sur NumPy, calcule les prédictions d'ensemble et génère un fichier téléchargeable enrichi des prédictions.
    \item \textbf{Diagnostic (\texttt{/analytics/}) } : Page affichant les courbes de performance pré-générées pour assurer la traçabilité des décisions des modèles.
\end{itemize}

\subsection{Pipeline de Déploiement Serverless (Vercel)}
L'environnement sans état de Vercel impose des contraintes de taille de package (limitation à 15MB compressés). Pour y répondre, nous utilisons le gestionnaire d'actifs statiques \textbf{WhiteNoise} qui intercepte et sert les fichiers multimédias et les CSS directement à travers le script de liaison WSGI de Django, configuré via le fichier \texttt{vercel.json} :
\begin{lstlisting}[language=json, caption=Configuration vercel.json de l'application Bank]
{
  "builds": [
    {
      "src": "bank_project/wsgi.py",
      "use": "@vercel/python",
      "config": {
        "maxLambdaSize": "15mb",
        "runtime": "python3.12"
      }
    }
  ],
  "routes": [
    {
      "src": "/(.*)",
      "dest": "bank_project/wsgi.py"
    }
  ]
}
\end{lstlisting}

% --- Section 7 ---
\section{Conclusion}
Ce projet démontre avec succès la pertinence opérationnelle d'un système d'apprentissage d'ensemble pour le scoring de prospects bancaires. En combinant des classificateurs aux comportements complémentaires (Régression Logistique, Forêt Aléatoire et Perceptron Multicouche), le modèle d'ensemble atteint un niveau de performance supérieur à celui de chaque modèle pris individuellement. Le déploiement serverless réalisé garantit l'accessibilité à grande échelle du moteur d'inférence avec des temps de réponse inférieurs à 100ms.

\newpage
\section*{Informations de Livraison}
\addcontentsline{toc}{section}{Informations de Livraison}
\begin{itemize}
    \item \textbf{Étudiant :} NEUSSI NJIETCHEU Patrice Eugene (Matricule : 21T2894)
    \item \textbf{Classe :} AIA4 (Intelligence Artificielle), ENSPY
    \item \textbf{Enseignant :} M. Bitha Junior
    \item \textbf{Dépôt de Code (GitHub) :} \url{https://github.com/neussi/bank_behavior_prediction}
    \item \textbf{Application en Ligne (Vercel) :} \url{https://bank-behavior-prediction.vercel.app}
\end{itemize}

\end{document}
